\section*{Аннотация}
\addcontentsline{toc}{section}{Аннотация}

Работа посвящена разработке программного комплекса для высокоэффективного решения задачи встречи космических аппаратов. Предложен подход, связывающий двухимпульсное решение в линейном приближении, нелинейное уточнение в возмущённой динамике и переход к манёврам конечной длительности. Состояние аппарата описывается модифицированными равноденственными элементами, что позволяет избежать особенностей кеплеровых элементов.

В линейной постановке двухимпульсная задача сведена к минимизации суммарной характеристической скорости по истинным долготам приложения импульсов. Для её решения применён метод Ньютона с множественным стартом; получены аналитические выражения для градиента и матрицы Гессе. Предложен способ дополнительного ускорения расчёта за счёт использования структуры локальных минимумов функции характеристической скорости.

Локально оптимальные решения линейной задачи используются для построения начального приближения в нелинейной многоимпульсной постановке. Уточнение выполняется в возмущённой модели движения при ограничениях на величину отдельных импульсов. Для учёта конечной длительности работы двигательной установки введена матрица чувствительности, описывающая линейное влияние манёвра с конечной тягой на изменение элементов орбиты.

Работоспособность алгоритмов подтверждена на задачах коррекции орбиты и фазирования. Результаты показывают, что разработанный комплекс позволяет быстро строить допустимые решения и уточнять их до достижения малой конечной невязки.
